\section*{الفصل \ref{ch:g9:natural-risks} --- العيش مع المخاطر الطبيعية}

\begin{solution}{exo:g9:natural-risks:1}
(أ) خطر. (ب) \omterm{def:g9:natural-risks:vulnerability}{هشاشة}. (ج) \omterm{def:g9:natural-risks:vulnerability}{هشاشة}. (د) خطر. (ه) كلاهما: فالسيول الطينية هي الخطر،
وبناء البلدة في طريقها يجعلها
هشة.
\end{solution}

\begin{solution}{exo:g9:natural-risks:2}
(أ) خطأ: حيث لا يسكن أحد، لا شيء يمكن أن يتضرر، \omterm{def:g9:natural-risks:risk}{فالمخاطر} منخفضة. (ب) صواب.
(ج) صواب. (د) خطأ: فالإنذار لا يغير \omterm{def:g4:volcanoes-earthquakes-first-look:volcano}{البركان}؛ بل يخفض \omterm{def:g9:natural-risks:vulnerability}{هشاشة} الناس، الذين
يستطيعون المغادرة في الوقت المناسب.
\end{solution}

\begin{solution}{exo:g9:natural-risks:3}
(أ) نعم، قبل أيام. (ب) نعم، قبل ساعات أو أيام. (ج) لا: \omterm{def:g4:volcanoes-earthquakes-first-look:earthquake}{فالزلازل} لا يمكن التنبؤ
بها. (د) نعم: فما إن يبدأ \omterm{def:g7:earthquakes:tsunami}{التسونامي}، حتى يستغرق ساعات لعبور \omterm{def:g2:water-on-earth:ocean}{المحيط}، والإنذار
ينتقل أسرع منه.
\end{solution}

\begin{solution}{exo:g9:natural-risks:4}
(أ) اهتزاز الأرض: أسراب من \omterm{def:g4:volcanoes-earthquakes-first-look:earthquake}{الزلازل} الصغيرة حين تصعد \omterm{def:g7:volcanoes:magma}{الصهارة}. (ب) ميل المنحدر:
انتفاخ \omterm{def:g4:volcanoes-earthquakes-first-look:volcano}{البركان}. (ج) كيف تتحرك الأرض، بالمليمترات: الانتفاخ حين تتجمع \omterm{def:g7:volcanoes:magma}{الصهارة}.
(د) غازات العمود الغازي: مزيد من الغاز ينبعث حين تصعد
\omterm{def:g7:volcanoes:magma}{الصهارة}.
\end{solution}

\begin{solution}{exo:g9:natural-risks:5}
(أ) $45 \div 2 = \qty{22.5}{km/h}$. (ب) $4 \times 2 = \qty{8}{km}$، وهي مسافة
كافية لبلوغ أرض أعلى. (ج) لم يرَ أحد \omterm{def:g4:volcanoes-earthquakes-first-look:eruption}{الثوران}، الذي حجبته \omterm{def:g5:natural-hazards:storm}{العاصفة}، ولم يُنذرهم أحد:
فلم تكن هناك وسيلة سريعة لإيصال المعلومات إلى البلدة، ولا خطة إخلاء. (د) مثلًا:
خطة إخلاء فيها طرق إلى أرض أعلى، يُتدرب عليها مسبقًا؛ وسلسلة إنذار من العلماء
إلى البلدة؛ ومزيد من الأجهزة على \omterm{def:g4:volcanoes-earthquakes-first-look:volcano}{البركان}.
\end{solution}

\begin{solution}{exo:g9:natural-risks:6}
(أ) \omterm{def:g5:water-shapes-the-land:flood}{الفيضان} الذي يعلو \qty{10}{m} فوق \omterm{def:g1:landscapes-around-us:water}{النهر}. (ب) نحو \qty{0.47}{km} \omterm{def:g5:water-shapes-the-land:flood}{لفيضان}
\qty{2}{m}، و\qty{0.56}{km} \omterm{def:g5:water-shapes-the-land:flood}{لفيضان} \qty{5}{m}، و\qty{0.67}{km} \omterm{def:g5:water-shapes-the-land:flood}{لفيضان}
\qty{10}{m} (والقراءات التي لا تبعد عنها أكثر من نحو \qty{0.05}{km} مقبولة).
(ج) \omterm{def:g5:water-shapes-the-land:flood}{فيضان} \qty{10}{m} وحده: فالأرض أعلى من 2 و\qty{5}{m} فوق \omterm{def:g1:landscapes-around-us:water}{النهر}. (د) يكفي
لبلوغه ارتفاع صغير \omterm{def:g1:landscapes-around-us:water}{للنهر}، \omterm{def:g5:water-shapes-the-land:flood}{والفيضانات} الصغيرة تأتي أكثر تكرارًا من الكبيرة.
\end{solution}

\begin{solution}{exo:g9:natural-risks:7}
(أ) خطر مرتفع \omterm{def:g9:natural-risks:vulnerability}{وهشاشة} منخفضة: \omterm{def:g9:natural-risks:risk}{مخاطر} متوسطة. (ب) خطر منخفض \omterm{def:g9:natural-risks:vulnerability}{وهشاشة} مرتفعة: \omterm{def:g9:natural-risks:risk}{مخاطر}
متوسطة. (ج) خطر مرتفع \omterm{def:g9:natural-risks:vulnerability}{وهشاشة} مرتفعة: \omterm{def:g9:natural-risks:risk}{مخاطر} عالية. (د) خطر منخفض \omterm{def:g9:natural-risks:vulnerability}{وهشاشة} منخفضة:
\omterm{def:g9:natural-risks:risk}{مخاطر} منخفضة.
\end{solution}

\begin{solution}{exo:g9:natural-risks:8}
(أ) C، ثم B، ثم A. (ب) تكوّن الدعامات مثلثات، والمثلث لا يمكن أن يميل دون أن
تتمدد دعامة أو تنضغط: فيحفظ الهيكل شكله بدل أن يميل. (ج) تلين العوازل حين تتحرك
الأرض ذهابًا وإيابًا تحتها، فلا يبلغ المدرسة إلا جزء أصغر بكثير من
الحركة.
\end{solution}

\begin{solution}{exo:g9:natural-risks:9}
(أ) $1 \div 100 = 1~\%$ في السنة. (ب) الفرصة هي نفسها كل سنة، مهما حدث في السنة
السابقة: فقد يأتي \omterm{def:g5:water-shapes-the-land:flood}{الفيضان} من جديد في السنة القادمة. (ج) $300 \div 100 = 3$
\omterm{def:g5:water-shapes-the-land:flood}{فيضانات}، في المتوسط.
\end{solution}

\begin{solution}{exo:g9:natural-risks:10}
(أ) $3600 \div 15 = 240$ قراءة. (ب) $240 \times 24 \times 365 =
\num{2102400} \approx 2.1 \times 10^6$ قراءة. (ج) $2400 \div 800 =
\qty{3}{h}$. (د) في هذه الساعات الثلاث، تستطيع مراكز الإنذار أن تنذر \omterm{def:g1:landscapes-around-us:water}{الساحل} ويستطيع
الناس الانتقال إلى أرض أعلى.
\end{solution}

\begin{solution}{exo:g9:natural-risks:11}
(أ) $4 \times 3 \times 3.8 = \qty{45.6}{L}$. (ب) $45.6 \div 1.5 = 30.4$،
أي 31 قارورة. (ج) مثلًا: طعام يبقى صالحًا، للأكل حين تُغلق المتاجر؛ ومصباح يدوي،
إذا انقطعت الكهرباء؛ ومذياع يعمل بالبطاريات أو باليد، لسماع الإنذارات؛ وحقيبة
إسعافات أولية، لعلاج الإصابات الصغيرة؛ وبطانيات دافئة. (د) كي يمكن أخذها فورًا إذا
اضطرت الأسرة إلى المغادرة بسرعة.
\end{solution}

\begin{solution}{exo:g9:natural-risks:12}
مثلًا: (1) رسم \omterm{def:g2:drawing-to-map:map}{خريطة} أخطار للسيول الطينية، \omterm{def:g1:landscapes-around-us:landform}{واديًا} بعد \omterm{def:g1:landscapes-around-us:landform}{وادٍ} (جانب الخطر)؛ (2) تركيب
\omterm{def:g7:earthquakes:seismometer}{مقاييس الزلازل} ومقاييس الميل ومستقبلات GPS وأجهزة قياس الغازات على \omterm{def:g4:volcanoes-earthquakes-first-look:volcano}{البركان} (جانب
الخطر)؛ (3) إقامة سلسلة إنذار سريعة من المرصد إلى المدينة، مع صفارات إنذار (جانب
\omterm{def:g9:natural-risks:vulnerability}{الهشاشة})؛ (4) تخطيط طرق إلى أرض أعلى ووضع إشارات عليها، والتدرب عليها بتمارين
(جانب \omterm{def:g9:natural-risks:vulnerability}{الهشاشة})؛ (5) عدم بناء بيوت أو مدارس جديدة في طريق السيول الطينية (جانب
\omterm{def:g9:natural-risks:vulnerability}{الهشاشة}).
\end{solution}

\begin{solution}{pb:g9:natural-risks:1}
\textbf{1.} يقع \omterm{def:g1:landscapes-around-us:water}{النهر} على \omterm{def:g4:world-physical-map:altitude}{ارتفاع} \qty{143}{m} بجانب C، و\qty{137}{m} بجانب A،
و\qty{126}{m} بجانب B؛ وعلى \omterm{def:g2:drawing-to-map:map}{الخريطة} يقع C في الغرب وB في الشرق، \omterm{def:g1:landscapes-around-us:water}{فالنهر} يجري
شرقًا.

\textbf{2.} A: $137 - 137 = \qty{0}{m}$; B: $276 - 126 = \qty{150}{m}$;
C: $236 - 143 = \qty{93}{m}$.

\textbf{3.} A وحده، فهو يرتفع أقل من \qty{5}{m} فوق \omterm{def:g1:landscapes-around-us:water}{النهر}.

\textbf{4.} نحو \qty{0.57}{km} (بين 0.5 و\qty{0.65}{km}).

\textbf{5.} الارتفاع على مسافة \qty{100}{m} هو نصف الارتفاع على مسافة
\qty{200}{m}: A $1.4 \div 2 = 0.7~\%$؛ B $83 \div 2 = 41.5~\%$؛
C $20 \div 2 = 10~\%$.

\textbf{6.} B: فانحداره شديد جدًا، \omterm{def:g3:stones-pebbles-sand:rock}{والصخر} المفكك والتربة قد ينزلقان على منحدر
شديد، ولا سيما بعد \omterm{def:g4:clouds-rain-snow:precipitation}{مطر} غزير.

\textbf{7.} A: نحو \qty{1}{km}؛ B: يكاد يقع على \omterm{def:g7:earthquakes:fault}{الصدع} (أقل من
\qty{0.05}{km})؛ C: نحو \qty{1.3}{km}.

\textbf{8.} إذا تحرك \omterm{def:g7:earthquakes:fault}{الصدع}، تنكسر الأرض وتنزاح على طوله، وهذا قد يمزق المبنى؛
\omterm{def:g4:volcanoes-earthquakes-first-look:earthquake}{وزلزال} على \omterm{def:g7:earthquakes:fault}{الصدع} سيهز الأرض هناك بشدة
بالغة.

\textbf{9.}

\begin{center}
{\small
\begin{tabular}{|c|c|c|c|}
\hline
الموقع & في نطاق \omterm{def:g5:water-shapes-the-land:flood}{الفيضان}؟ & منحدر شديد؟ & على \omterm{def:g7:earthquakes:fault}{الصدع}؟ \\
\hline
A & نعم & لا & لا \\
B & لا & نعم & نعم \\
C & لا & لا & لا \\
\hline
\end{tabular}}
\end{center}

\textbf{10.} الموقع C: فهو الموقع الوحيد الذي حصل على ثلاث «لا»: \qty{93}{m} فوق
\omterm{def:g1:landscapes-around-us:water}{النهر}، وانحدار لطيف نسبته \qty{10}{\%}، وبُعد يزيد على \qty{1}{km}
عن \omterm{def:g7:earthquakes:fault}{الصدع}.

\textbf{11.} مثلًا: دعامات قطرية وجدران قص في الهيكل، ووصلات متينة، أو عوازل تحت
المبنى؛ ويتدرب التلاميذ على تمارين \omterm{def:g4:volcanoes-earthquakes-first-look:earthquake}{الزلازل} (انبطح، واحتمِ،
وتمسّك).

\textbf{12.} \textbf{\babelsublr{17} بيتًا} من 28 تقع داخل نطاق \omterm{def:g5:water-shapes-the-land:flood}{الفيضان}:
$17 \div 28 \approx 0.61$، أي نحو 3 بيوت من كل 5.
\end{solution}
